\documentclass[11pt,a4paper]{article}

\usepackage[utf8]{inputenc}
\usepackage[margin=0.85in]{geometry}
\usepackage{amsmath,amsfonts,amssymb,mathtools}
\usepackage{booktabs,array,xcolor,hyperref}
\usepackage{listings,tcolorbox}
\usepackage{enumitem,titlesec}
\usepackage{microtype}

% Color definitions
\definecolor{primarynavy}{HTML}{0A192F}
\definecolor{accentblue}{HTML}{00B4D8}
\definecolor{darkslate}{HTML}{1E293B}
\definecolor{cardbg}{HTML}{F8FAFC}
\definecolor{cardborder}{HTML}{CBD5E1}
\definecolor{answertitle}{HTML}{166534}
\definecolor{answerbox}{HTML}{F0FDF4}
\definecolor{answerborder}{HTML}{BBF7D0}

\hypersetup{
    colorlinks=true,
    linkcolor=accentblue,
    urlcolor=accentblue,
    citecolor=primarynavy,
    pdftitle={Applied Quantitative Research: Course Syllabus & Technical Leveling Guide},
    pdfauthor={Lucca Simeoni Pavan, Ph.D.}
}

% Section styling
\titleformat{\section}{\color{primarynavy}\Large\bfseries}{\thesection.}{0.5em}{}[\titlerule]
\titleformat{\subsection}{\color{primarynavy}\large\bfseries}{\thesubsection}{0.5em}{}
\titleformat{\subsubsection}{\color{darkslate}\normalsize\bfseries}{\thesubsubsection}{0.5em}{}

% Listings styling for Python
\lstdefinestyle{pythonstyle}{
    language=Python,
    basicstyle=\ttfamily\footnotesize,
    keywordstyle=\color{primarynavy}\bfseries,
    commentstyle=\color{gray}\itshape,
    stringstyle=\color{accentblue},
    backgroundcolor=\color{cardbg},
    frame=single,
    rulecolor=\color{cardborder},
    breaklines=true,
    showstringspaces=false,
    numbers=left,
    numberstyle=\tiny\color{gray},
    numbersep=8pt,
    tabsize=4
}
\lstset{style=pythonstyle}

% Box definitions
\newtcolorbox{infobox}[1][]{
    colback=cardbg,
    colframe=primarynavy,
    fonttitle=\bfseries\color{white},
    title=#1,
    arc=3mm,
    boxrule=1pt,
    left=10pt, right=10pt, top=8pt, bottom=8pt
}

\newtcolorbox{answerbox}[1][]{
    colback=answerbox,
    colframe=answerborder,
    coltitle=answertitle,
    title=#1,
    arc=2mm,
    boxrule=1pt,
    left=8pt, right=8pt, top=6pt, bottom=6pt
}

\begin{document}

% ----------------- TITLE BLOCK -----------------
\begin{center}
    {\Huge \textbf{\color{primarynavy} APPLIED QUANTITATIVE RESEARCH}}\\[6pt]
    {\Large \textbf{\color{accentblue} Systematic Modeling, Multi-Factor Investing \& Risk Allocation in Python}}\\[10pt]
    \rule{\textwidth}{1.5pt}\\[8pt]
    {\normalsize \textbf{Lucca Simeoni Pavan, Ph.D.}}\\
    {\small Former Head of Quantitative Strategies \& Portfolio Allocation Manager $\bullet$ Ph.D. in Economics}\\
    {\small \textit{Official Course Curriculum, Technical Leveling Guide \& Diagnostic Assessment (30h Institutional Formation)}}\\[4pt]
    \rule{\textwidth}{0.8pt}
\end{center}

\vspace{10pt}

% ----------------- SECTION 1: SYLLABUS -----------------
\section{Official Course Curriculum \& Syllabus}

\begin{infobox}[Program Overview]
\textbf{Program Name:} Applied Quantitative Research: Systematic Modeling, Factor Investing \& Risk Management in Python\\
\textbf{Workload:} 30 Hours (Live Sessions + Recorded Architecture Reviews + Coding Laboratories)\\
\textbf{Tech Stack:} Python 3.10+, Pandas, NumPy, Statsmodels, SciPy, Matplotlib, Seaborn, yfinance\\
\textbf{Core Objective:} Empower researchers, economists, data scientists, and investment analysts to architect end-to-end institutional quantitative pipelines---from point-in-time financial data hygiene to bias-free walk-forward backtesting and Ledoit-Wolf / HRP risk parity portfolio optimization.
\end{infobox}

\subsection*{Module 1: Financial Data Pipelines \& Point-in-Time Hygiene (5 Hours)}
\begin{itemize}[leftmargin=1.5em]
    \item \textbf{Lecture 1.1:} Pipeline Architecture: Multidimensional time-series structures (multi-index DataFrames and panels).
    \item \textbf{Lecture 1.2:} Institutional Data Ingestion: Automated feeds from central bank APIs, SEC/CVM filings, and exchanges.
    \item \textbf{Lecture 1.3:} Corporate Actions: Cash dividends, splits, reverse splits, spin-offs, and ANBIMA/NYSE calendar conventions.
    \item \textbf{Lecture 1.4:} Eliminating Modeling Biases: Eradicating survivorship bias and look-ahead bias via Point-in-Time (PIT) lags.
    \item \textbf{Hands-on Laboratory:} Building an auditable, survivorship-free 10-year historical equity data universe.
\end{itemize}

\subsection*{Module 2: Return Mathematics \& Institutional Risk Metrics (5 Hours)}
\begin{itemize}[leftmargin=1.5em]
    \item \textbf{Lecture 2.1:} Simple vs. Logarithmic Returns: Algebraic properties, cross-sectional vs. temporal additivity.
    \item \textbf{Lecture 2.2:} Empirical Return Distributions: Heavy tails, negative skewness, excess kurtosis, and non-normality.
    \item \textbf{Lecture 2.3:} Asymmetric Metrics: Sharpe Ratio, Sortino Ratio (downside semi-variance), Calmar Ratio, Information Ratio.
    \item \textbf{Lecture 2.4:} Maximum Drawdown Dynamics: Analytical peak-to-trough computation, drawdown duration, and recovery paths.
    \item \textbf{Lecture 2.5:} Tail Risk: Parametric/Historical Value at Risk ($\text{VaR}_{95\%}$) and Conditional VaR ($\text{CVaR}$ / Expected Shortfall).
    \item \textbf{Hands-on Laboratory:} Automated pipeline generating institutional Risk Tear Sheets from any input price matrix.
\end{itemize}

\subsection*{Module 3: Multi-Factor Investing \& Alpha Engineering (5 Hours)}
\begin{itemize}[leftmargin=1.5em]
    \item \textbf{Lecture 3.1:} Evolution of Asset Pricing: From Markowitz-Sharpe Single-Index CAPM to Fama-French 5 Factors and Carhart.
    \item \textbf{Lecture 3.2:} Momentum Anomalies: Cross-Sectional Momentum (winners vs. losers) vs. Time-Series Momentum (trend-following).
    \item \textbf{Lecture 3.3:} Fundamental Factors: Value ($E/P$, $B/M$, $EV/EBITDA$) and Quality ($ROE$, $ROIC$, Operating Margin, Accruals).
    \item \textbf{Lecture 3.4:} Risk Premia: Low Beta / Low Volatility anomaly and Size (Small-Cap premium).
    \item \textbf{Lecture 3.5:} Surviving the ``Factor Zoo'': Multiple hypothesis testing, Harvey-Liu-Bikker $t > 3.0$ thresholds, Bonferroni adjustments.
    \item \textbf{Hands-on Laboratory:} Multi-factor scoring engine with cross-sectional Z-Scores and Frisch-Waugh-Lovell (FWL) orthogonalization.
\end{itemize}

\subsection*{Module 4: Realistic Institutional Backtesting Frameworks (5 Hours)}
\begin{itemize}[leftmargin=1.5em]
    \item \textbf{Lecture 4.1:} Backtest Paradigms: Vectorized (rapid research screening) vs. Event-Driven (microstructure and execution fidelity).
    \item \textbf{Lecture 4.2:} Friction Accounting: Brokerage commissions, exchange fees, settlement costs, short-borrow rates, and taxes.
    \item \textbf{Lecture 4.3:} Slippage \& Market Impact: Average Daily Trading Volume (ADTV) participation bounds and Almgren-Chriss square-root impact.
    \item \textbf{Lecture 4.4:} Dynamic Rebalancing: Turnover constraints, buffer zones, transaction fee drag, and optimal trade scheduling.
    \item \textbf{Lecture 4.5:} Cross-Validation: Purged \& Embargoed K-Fold cross-validation (de Prado) avoiding leakage in serial correlation.
    \item \textbf{Hands-on Laboratory:} Full 5-year multi-factor equity strategy backtest under institutional execution constraints.
\end{itemize}

\subsection*{Module 5: Portfolio Optimization \& Risk Allocation (5 Hours)}
\begin{itemize}[leftmargin=1.5em]
    \item \textbf{Lecture 5.1:} Markowitz Mean-Variance Optimization: Sample covariance instability and the ``error maximization'' trap.
    \item \textbf{Lecture 5.2:} Covariance Regularization: Ledoit-Wolf analytical shrinkage and Random Matrix Theory (RMT) eigenvalue cleaning.
    \item \textbf{Lecture 5.3:} Risk Parity \& Equal Risk Contribution: Decomposing marginal risk contributions across volatile asset baskets.
    \item \textbf{Lecture 5.4:} Hierarchical Risk Parity (HRP): Unsupervised machine learning dendrogram clustering bypassing matrix inversion.
    \item \textbf{Hands-on Laboratory:} Empirical horse-race benchmark: $1/N$ Equal Weight vs. Markowitz vs. Ledoit-Wolf vs. HRP.
\end{itemize}

\subsection*{Module 6: Capstone Project \& Production Quant Deployment (5 Hours)}
\begin{itemize}[leftmargin=1.5em]
    \item \textbf{Lecture 6.1:} Production Code Architecture: Modular classes, static typing, unit testing, and reproducible research notebooks.
    \item \textbf{Lecture 6.2:} Automated Reporting: Compiling interactive HTML/PDF institutional risk sheets for investment committees.
    \item \textbf{Capstone Challenge:} End-to-end development of a proprietary systematic trading strategy with 1-on-1 code review.
\end{itemize}

\vspace{10pt}

% ----------------- SECTION 2: TECHNICAL LEVELING GUIDE -----------------
\section{Technical Leveling Guide \& Mathematical Formulations}

\subsection{Simple Returns vs. Logarithmic Returns}
A foundational axiom in quantitative finance is the strict separation between arithmetic (simple) and logarithmic returns:
\begin{align}
R_t &= \frac{P_t - P_{t-1}}{P_{t-1}} = \frac{P_t}{P_{t-1}} - 1 \quad \text{(Simple Arithmetic Return)} \\
r_t &= \ln\left(\frac{P_t}{P_{t-1}}\right) = \ln(P_t) - \ln(P_{t-1}) \quad \text{(Logarithmic Continuously-Compounded Return)}
\end{align}

\subsubsection*{The Golden Rules of Quantitative Return Aggregation}
\begin{enumerate}
    \item \textbf{Cross-Sectional Additivity (Across Assets):} Simple returns aggregate linearly across assets in a portfolio weighted by $w_i \in \mathbb{R}^N$:
    \begin{equation}
    R_{p,t} = \sum_{i=1}^N w_i R_{i,t} \quad \text{where } \sum_{i=1}^N w_i = 1
    \end{equation}
    \textit{Never average or sum weighted log returns to compute portfolio value!}
    \item \textbf{Temporal Additivity (Across Time):} Log returns aggregate linearly across time intervals. The total compound return over $T$ periods is simply:
    \begin{equation}
    r_{\text{total}} = \sum_{t=1}^T r_t \implies R_{\text{total}} = \exp\left(\sum_{t=1}^T r_t\right) - 1
    \end{equation}
    Use log returns for time-series econometric regressions, volatility modeling (GARCH), and Brownian motion assumptions.
\end{enumerate}

\subsection{Annualization Conventions (252 Business Days)}
Under the assumption of independent and identically distributed ($i.i.d.$) returns:
\begin{equation}
\bar{R}_{\text{annual}} = (1 + \bar{R}_{\text{daily}})^{252} - 1, \qquad \sigma_{\text{annual}} = \sigma_{\text{daily}} \times \sqrt{252}
\end{equation}

\subsection{Institutional Performance \& Tail Risk Metrics}
\begin{itemize}[leftmargin=1.5em]
    \item \textbf{Sharpe Ratio:} Measures excess return over the risk-free rate ($R_f$) per unit of total standard deviation:
    \begin{equation}
    \text{Sharpe} = \frac{\bar{R}_p - R_f}{\sigma_p}
    \end{equation}
    \item \textbf{Sortino Ratio:} Replaces total volatility with downside semi-deviation ($\sigma_{\text{down}}$), isolating adverse volatility:
    \begin{equation}
    \text{Sortino} = \frac{\bar{R}_p - R_f}{\sigma_{\text{down}}}, \quad \text{where } \sigma_{\text{down}} = \sqrt{\frac{1}{T}\sum_{t=1}^T \min(0, R_{p,t} - R_f)^2} \times \sqrt{252}
    \end{equation}
    \item \textbf{Maximum Drawdown (MDD):} The greatest historical peak-to-trough capital decline:
    \begin{equation}
    \text{MDD} = \min_{t \in [0, T]} \left( \frac{\text{NAV}_t - \max_{\tau \le t} \text{NAV}_\tau}{\max_{\tau \le t} \text{NAV}_\tau} \right)
    \end{equation}
    \item \textbf{Conditional Value at Risk ($\text{CVaR}_\alpha$ / Expected Shortfall):} Expected loss conditional on exceeding the $\text{VaR}_\alpha$ threshold:
    \begin{equation}
    \text{CVaR}_\alpha = \mathbb{E}\left[ L \mid L \ge \text{VaR}_\alpha \right] = \frac{1}{1-\alpha} \int_\alpha^1 \text{VaR}_u \, du
    \end{equation}
\end{itemize}

\subsection{Factor Orthogonalization (Frisch-Waugh-Lovell Theorem)}
To isolate pure idiosyncratic alpha from known systematic factors (e.g. Market Beta, Size, Value):
\begin{equation}
M_X = I - X(X^T X)^{-1} X^T \implies \tilde{y} = M_X y, \quad \tilde{Z} = M_X Z \implies \hat{\gamma} = (\tilde{Z}^T \tilde{Z})^{-1} \tilde{Z}^T \tilde{y}
\end{equation}

\subsection{Ledoit-Wolf Covariance Shrinkage}
Shrinking an empirical sample covariance matrix $S$ toward a highly structured shrinkage target $F$ (constant correlation matrix) with optimal shrinkage intensity $\delta^* \in [0, 1]$:
\begin{equation}
\hat{\Sigma}_{\text{LW}} = \delta^* F + (1 - \delta^*) S
\end{equation}

\subsection{Python Production Scorecard Engine}
The following script executes institutional return and tail-risk calculations:

\begin{lstlisting}
import numpy as np
import pandas as pd
import yfinance as yf

# 1. Download adjusted closing prices
tickers = ['SPY', 'QQQ', 'AAPL', 'NVDA', 'TLT']
data = yf.download(tickers, start='2021-01-01', end='2026-01-01', progress=False)
prices = data['Adj Close'].dropna()
returns = prices.pct_change().dropna()

# 2. Institutional Quantitative Risk Function
def compute_scorecard(series: pd.Series, rf_annual: float = 0.045) -> pd.Series:
    rf_daily = (1 + rf_annual) ** (1 / 252) - 1
    ann_return = (1 + series.mean()) ** 252 - 1
    ann_vol = series.std() * np.sqrt(252)
    sharpe = (ann_return - rf_annual) / ann_vol if ann_vol > 0 else 0.0
    
    excess_ret = series - rf_daily
    downside_vol = excess_ret[excess_ret < 0].std() * np.sqrt(252)
    sortino = (ann_return - rf_annual) / downside_vol if downside_vol > 0 else 0.0
    
    cum_ret = (1 + series).cumprod()
    mdd = ((cum_ret - cum_ret.cummax()) / cum_ret.cummax()).min()
    
    return pd.Series({
        'Ann. Return': f"{ann_return * 100:.2f}%",
        'Ann. Volatility': f"{ann_vol * 100:.2f}%",
        'Sharpe Ratio': f"{sharpe:.2f}",
        'Sortino Ratio': f"{sortino:.2f}",
        'Max Drawdown': f"{mdd * 100:.2f}%"
    })

print(returns.apply(compute_scorecard).to_string())
\end{lstlisting}

\vspace{10pt}

% ----------------- SECTION 3: DIAGNOSTIC TEST -----------------
\section{Diagnostic Assessment: 10 Institutional Questions}

\begin{infobox}[Question 1 $\bullet$ Return Mathematics]
When constructing a portfolio allocating 40\% in Asset A and 60\% in Asset B, which formulation correctly computes the daily total return?\\
A) The weighted arithmetic average of the logarithmic returns.\\
\textbf{B) The weighted arithmetic average of the simple arithmetic returns: $R_{p,t} = \sum w_i R_{i,t}$.}\\
C) The geometric product of simple returns.\\
D) The logarithm of closing prices divided by volume.
\end{infobox}

\begin{infobox}[Question 2 $\bullet$ Backtest Biases \& Microstructure]
A researcher builds a multi-factor strategy using Q4 financial statements (fiscal period ending Dec 31) to rebalance stock holdings on the first trading session of January. What fatal quantitative error has occurred?\\
A) Survivorship bias due to delisted companies.\\
\textbf{B) Look-ahead bias, because audited annual 10-K financial statements are publicized months later in March or April.}\\
C) Specification error by using calendar days instead of trading days.\\
D) Convexity bias in volatility annualization.
\end{infobox}

\begin{infobox}[Question 3 $\bullet$ Python Vectorization]
Given a date-indexed DataFrame \texttt{df} containing adjusted prices in column \texttt{'close'}, which statement correctly computes daily simple arithmetic returns?\\
A) \texttt{df['close'].diff() / df['close']}\\
B) \texttt{np.log(df['close']) - np.log(df['close'].shift(-1))}\\
\textbf{C) \texttt{df['close'].pct_change().dropna()}}\\
D) \texttt{df['close'].rolling(252).mean()}
\end{infobox}

\begin{infobox}[Question 4 $\bullet$ Volatility Annualization Standards]
An equity asset displays a daily return standard deviation of 2.0\% in a market with 252 business days. Under the $i.i.d.$ assumption, what is its annualized volatility?\\
A) $2.0\% \times 252 = 504.0\%$\\
\textbf{B) $2.0\% \times \sqrt{252} \approx 31.75\%$}\\
C) $2.0\% / \sqrt{252} \approx 0.126\%$\\
D) $(1 + 0.02)^{252} - 1 \approx 145.2\%$
\end{infobox}

\begin{infobox}[Question 5 $\bullet$ Asymmetric Risk Metrics]
Why do quantitative managers favor the Sortino Ratio over the Sharpe Ratio for strategies displaying positive return skewness?\\
A) The Sortino Ratio ignores the risk-free rate entirely.\\
\textbf{B) The Sharpe Ratio penalizes upside return spikes identically to drawdowns, whereas Sortino penalizes only downside volatility.}\\
C) The Sortino Ratio is independent of sample size.\\
D) The Sharpe Ratio cannot be calculated when returns are non-negative.
\end{infobox}

\begin{infobox}[Question 6 $\bullet$ Cross-Validation in Financial Time Series]
Why is standard Scikit-Learn K-Fold cross-validation invalid and dangerous for financial asset pricing models?\\
\textbf{A) Standard K-Fold shuffles observations and trains on future data to predict the past, causing temporal data leakage.}\\
B) K-Fold is computationally restricted to classification tasks.\\
C) The number of folds must match the number of assets.\\
D) K-Fold artificially doubles transaction fee estimates.
\end{infobox}

\begin{infobox}[Question 7 $\bullet$ Factor Modeling Paradigms]
What is the core distinction between Cross-Sectional Momentum and Time-Series Momentum?\\
A) Cross-Sectional is applied to fixed income and Time-Series is applied to equities.\\
\textbf{B) Cross-Sectional ranks assets relative to peers at a timestamp (buying top decile, shorting bottom decile), while Time-Series evaluates each asset's own historical trend in absolute terms.}\\
C) Cross-Sectional uses moving averages while Time-Series uses regressions.\\
D) Both terms are synonymous in empirical asset pricing literature.
\end{infobox}

\begin{infobox}[Question 8 $\bullet$ Microstructure Frictions \& Execution]
When backtesting a systematic small-cap strategy with \$20M in AUM, which parameter is crucial to prevent illusory alpha?\\
A) CPI inflation lags.\\
\textbf{B) Average Daily Trading Volume (ADTV) limits (e.g. max 5--10\% ADTV) and nonlinear market impact models.}\\
C) Social media follower count of listed management.\\
D) Book-to-market ratio divided by shares outstanding.
\end{infobox}

\begin{infobox}[Question 9 $\bullet$ Covariance Estimation \& Markowitz]
What is the central empirical drawback of traditional Markowitz Mean-Variance Optimization applied to sample covariance matrices?\\
A) The algorithm cannot process positive returns.\\
\textbf{B) Inverting an ill-conditioned sample covariance matrix acts as an ``error maximizer'', concentrating extreme unstable weights on noisy assets; solved by Ledoit-Wolf shrinkage.}\\
C) Markowitz requires that all portfolio returns be zero.\\
D) The solver requires unconstrained shorting.
\end{infobox}

\begin{infobox}[Question 10 $\bullet$ Hierarchical Risk Parity (HRP)]
What breakthrough did Marcos L\'opez de Prado's Hierarchical Risk Parity (HRP) algorithm introduce?\\
\textbf{A) It uses machine learning tree clustering on the correlation matrix to allocate risk recursively, completely bypassing covariance matrix inversion.}\\
B) It guarantees 0\% drawdown across all market regimes.\\
C) It replaces mathematical solvers with NLP models.\\
D) It optimizes option strike prices for high-frequency trading.
\end{infobox}

\vspace{10pt}

% ----------------- SECTION 4: DIAGNOSTIC RUBRIC -----------------
\section{Diagnostic Rubric \& Program Fit}

\begin{table}[h]
\centering
\begin{tabular}{p{0.2\textwidth} p{0.35\textwidth} p{0.38\textwidth}}
\toprule
\textbf{Score} & \textbf{Institutional Diagnostic Level} & \textbf{Recommended Program Focus} \\
\midrule
\textbf{0 to 4 Correct} & \textbf{Level 1: Foundations / In Transition}\newline Analytical background, but generic habits from classical data science. & \textbf{Modules 1 \& 2} are critical for building point-in-time data hygiene, clean series, and tail risk metrics. \\
\midrule
\textbf{5 to 7 Correct} & \textbf{Level 2: Intermediate Quant Analyst}\newline Solid Python syntax and basic stats, with gaps in friction modeling and multi-factor isolation. & \textbf{Modules 3 \& 4} will elevate backtesting to buy-side desk standards (FWL partialing, slippage, and factor spreads). \\
\midrule
\textbf{8 to 10 Correct} & \textbf{Level 3: Quant Desk Ready}\newline Strong conceptual maturity and sharpened intuition for microstructure and risk budgeting. & \textbf{Modules 4, 5 \& 6} will solidify mastery of Ledoit-Wolf, HRP allocators, and production-grade take-home challenges. \\
\bottomrule
\end{tabular}
\end{table}

\vspace{10pt}

\begin{answerbox}[Exclusive VIP Waitlist Advantage]
By registering your email on the official landing page, you have secured your priority spot for the 1st Cohort of the \textbf{Applied Quantitative Research Program} with an \textbf{exclusive 20\% launch discount} and direct code review mentoring sessions with Dr. Lucca Simeoni Pavan.
\end{answerbox}

\end{document}
